\section{Sân tập}
\label{sec:san-tap}

\textbf{Sân tập} là một dự án \emph{riêng tư} dành cho việc làm bài tập. Không
ai khác có thể xem dự án này, nên bạn cứ thoải mái thử nghiệm và chỉnh sửa tùy
ý.

\subsection{Mở sân tập}
\begin{enumerate}
  \item Mở \textbf{Sổ tay SumiTeX} (có sẵn trong mục \nut{Shared with you},
    hoặc vào menu \nut{Help} → \nut{SumiTeX guide}).
  \item Nhấp vào nút \nut{Open my playground} trên thanh thông báo nhỏ ở đầu
    trang.
  \item Ở lần đầu tiên, SumiTeX sẽ tạo cho bạn một bản sao riêng mang tên
    \textbf{Sân tập SumiTeX}. Từ những lần sau, nút này sẽ mở lại đúng sân tập
    đó.
\end{enumerate}

\subsection{Bên trong có gì}
Mỗi bài tập tương ứng với một tệp từ \texttt{bai-01.tex} đến
\texttt{bai-06.tex}, trong đó từng yêu cầu cần làm được đánh dấu bằng lệnh
\lenh{todo\{…\}}:
\begin{itemize}
  \item Trên bản PDF, các việc chưa hoàn thành sẽ hiển thị dưới dạng ghi chú ở
    lề trang.
  \item Trong bảng \textbf{Mục lục} (khi bật \textbf{Blocks}), mỗi yêu cầu của
    tệp đang mở sẽ hiện thành một dòng \emph{ghi chú}; chẳng hạn, khi mở
    \texttt{bai-02.tex}, bạn sẽ thấy danh sách việc của bài 2. Sau khi làm
    xong, bạn chỉ cần bấm ✎ rồi chọn \nut{✓ Done} để ẩn ghi chú đó đi.
  \item Nếu gặp khó khăn, bạn có thể tham khảo lời giải tại
    phụ lục~\ref{sec:loi-giai}.
\end{itemize}

\begin{meo}
Nếu muốn làm lại từ đầu, bạn hãy bấm nút \nut{Start over} trên thanh công cụ
phía trên của sân tập. Bản cũ vẫn sẽ được giữ lại với hậu tố ``(cũ)'' để bạn
tiện đối chiếu hoặc tự xóa sau.
\end{meo}

\begin{chuy}
Khi có bài tập mới được cập nhật, thanh trên cùng của sân tập sẽ hiển thị thông
báo \nut{new exercises are ready}. Bạn hãy bấm \nut{Get the new exercises} để
tạo một sân tập mới; toàn bộ bài bạn đã làm trước đó vẫn được bảo toàn nguyên
vẹn trong sân tập cũ.
\end{chuy}